
\subsubsection{Generalization Gap in Deep Learning}

The disconnect between benchmark and deployment performance has been documented across multiple studies. Recht et al. (2019) created ImageNetV2 by replicating the original data collection process and found that all tested classifiers experienced 11-14\% accuracy drops. These studies measure gaps through evaluation metrics but do not identify model-internal signals that might predict such gaps.

\subsubsection{Underspecification and Shortcut Learning}

D'Amour et al. (2020) formalized underspecification as a challenge: multiple models can achieve equivalent training performance while exhibiting different behaviors under distribution shift. Wang et al. (2025) showed that ImageNet-trained CNNs learn frequency shortcuts. Hermann et al. (2020) and Geirhos et al. (2019) documented texture bias in convolutional networks. Our benchmark fingerprint framework generalizes this analysis, asking whether fine-tuning creates detectable signatures regardless of the specific spurious mechanism.

\subsubsection{Linear Probing for Representation Analysis}

Linear probing has become a standard tool for understanding what information neural network representations encode (Alain \& Bengio, 2017; Kornblith et al., 2019). We apply this methodology to detect benchmark identity from fine-tuned representations.

\subsubsection{Our Positioning}

Prior work establishes that generalization gaps exist, underspecification causes divergence, and shortcuts can be domain-specific. We contribute a representation-level perspective: measuring benchmark fingerprints directly using linear probing. Our negative finding---that BFS does not correlate with gap under tested conditions---identifies limitations of the proposed metric and opens directions for alternative approaches.

